% Complete candidate proof in the explicitly stated constant-case scope.
\begin{proof}[Proof of \cref{prop:p081-uniform-fibre-fs-extension}]
Put \(R=k[[t]]\). First use the constant-model reference.

\emph{Integral sections detected by one generic Gauss point.}
Let \(V=H^0(X,mL)\). Flat base change gives
\(H^0(X_K,mL_K)=V\otimes_k K\): compute global sections as the
equalizer for a finite affine cover and tensor with the flat field
extension. Projectivity makes \(V\) finite dimensional. The generating
system over \(K\) implies that \(mL\) is globally generated: the
evaluation cokernel vanishes after faithful flat base change. Choose a
generating \(k\)-system \((q_\ell)\) for it.

Every \(k\)-section has trivial norm at most one, and
\(\max_\ell|q_\ell|_0=1\). These assertions also hold on \(F\) for
the constant-model metric. Indeed, take the centre of an analytic point
on the proper \(R\)-model and a local frame there. All section
coefficients are integral; one of the generating coefficients is a unit.
This works at points with proper support as well.

Let \(\xi\in F\) be the Gauss extension of the generic trivial
valuation on \(X\). For a \(k\)-basis \((e_\ell)\) of \(V\),
\begin{equation}\label{eq:p081-uniform-integral-norm}
 \left|\sum_\ell a_\ell(t)e_\ell\right|_0(\xi)
          =\max_\ell|a_\ell(t)|,\qquad a_\ell(t)\in K.
\end{equation}
To check this, expand in Laurent powers of \(t\) in a rational frame.
The first nonzero coefficient is a nonzero \(k\)-linear combination of
the \(e_\ell\), hence has generic trivial valuation zero. The Gauss
formula \cite[\S1.3]{BJ22} evaluates the section at its first exponent.
This uses only that formula, not the erroneous full-image identification
printed in Remark~1.9. A section has norm at most one at \(\xi\) if
and only if it belongs to \(V\otimes R\).

\emph{Normalize weights, then truncate with a strict margin.}
In \eqref{eq:p081-uniform-fs-definition} write
\(\lambda_j=n_j+\alpha_j\), with \(n_j\in\mathbf Z\) and
\(0\leq\alpha_j<1\), and put \(a_j=t^{-n_j}s_j\). Then
\[
                  \mathrm e^{mf}=\max_j
                       \mathrm e^{\alpha_j}|a_j|_0.
\]
The upper bound \(f\leq0\), evaluated at \(\xi\), implies
\(a_j\in V\otimes R\) by \eqref{eq:p081-uniform-integral-norm}; its
order is in fact at least \(\lceil\alpha_j\rceil\). Write
\(a_j=\sum_{n\geq0}a_{j,n}t^n\), \(a_{j,n}\in V\). Set
\(D=mC\) and \(b_j=\sum_{n=0}^{D}a_{j,n}t^n\). On every point of
the full fibre,
\begin{equation}\label{eq:p081-uniform-tail}
 \mathrm e^{\alpha_j}|a_j-b_j|_0
 \leq\mathrm e^{\alpha_j-D-1}<\mathrm e^{-D}
 \leq\mathrm e^{-mM}\leq\mathrm e^{mf}.
\end{equation}
Each coefficient section has norm at most one; finite partial sums give
the bound by the ultrametric inequality, and the infinite sum converges
in the completed residue field. Alternatively expand in the finite basis
\((e_\ell)\); its coefficient tails converge uniformly with the same
bound. Thus the estimate is uniform and includes proper-support points.

If finitely many vectors \(a_j,b_j\) in a non-Archimedean normed line
satisfy \(w_j|a_j-b_j|<B=\max_j w_j|a_j|>0\), then
\(\max_j w_j|b_j|=B\). All terms are at most \(B\); at an index
attaining \(B\), a strictly smaller error cannot change the norm.
Apply this to \eqref{eq:p081-uniform-tail}. It proves the exact identity
\begin{equation}\label{eq:p081-uniform-truncation}
                f=\frac1m\max_j\{\log|b_j|_0+\alpha_j\}
                      \quad\hbox{on all of }F.
\end{equation}
Cancellation among retained terms at a general analytic point is allowed.
We never apply a Gauss no-cancellation rule at such a point; only the
strict tail margin is used there.

\emph{Homogenize and pad with negatively weighted generators.}
In \(H^0(S,mA)\) take
\[
 B_j=\sum_{n=0}^{D}a_{j,n}T^nU^{D-n},\qquad
 Q_{\ell,0}=q_\ell U^D,\qquad Q_{\ell,\infty}=q_\ell T^D.
\]
Give \(B_j\) weight \(\alpha_j\), and both padding systems weight
\(-mC\). The padding sections alone have no common zero, including
over \(t=0\) and infinity. Their weighted maximum with the \(B_j\)
defines \(\Phi\in H_{\mathbf R}(A)\).

On \(\eta\), the trivial base metric has \(|U|_0=1\),
\(|T|_0=\mathrm e^{-1}\). In the frame \(U^D\), the first sections
restrict to \(b_j\), and each weighted padding term is at most
\(\mathrm e^{-mC}\leq\mathrm e^{mf}\). Hence
\eqref{eq:p081-uniform-truncation} gives \(\Phi|_F=f\).
Globally, all section norms are at most one, so
\(\Phi\leq1/m\leq1\). Also
\(\max_{\ell,\epsilon\in\{0,\infty\}}|Q_{\ell,\epsilon}|_0=1\),
because both original systems generate. Thus \(\Phi\geq-C\).
The normalized weights lie in \([-C,1)\), and rationality is preserved.
For rational weights one can first clear denominators by powers as in
\cite[(2.9)]{BJ22}; then all residual weights can be zero and the upper
bound improves to zero. This is optional.

Negative FS weights are essential: a polynomial over trivially valued
\(k\) with a nonzero constant term cannot itself be small at \(\eta\).
The weight supplies the smallness while the underlying section removes
the central-fibre base locus. No blowup or further uncontrolled twist is
being suppressed.

For the reference \(\rho\), let \(\psi=\rho_K+\phi\) be the
absolute FS potential relative to the constant model. Then
\(-(M+a+b)\leq\psi-a\leq0\). Extend \(\psi-a\) as above and put
\(H=\Phi+a-p^*\rho\). Its underlying metric has FS absolute potential
\(\Phi+a\); its restriction and stated bounds follow immediately.
The relative potential \(H\) is not being confused with an FS function
relative to the trivial reference.
\end{proof}

\begin{proof}[Proof of \cref{thm:p081-constant-laurent-envelope-transfer-candidate}]
\emph{A collar of the full analytic fibre.}
On the open set \(\mathcal C=\{0<|t|<1\}\subset S^{\mathrm{an}}\),
put \(a(y)=-\log|t(y)|>0\). Its base point is necessarily
\(\eta_{a(y)}=a(y)\operatorname{ord}_{t=0}\), since a base polynomial
over \(k\) has a unique lowest nonzero power and distinct powers have
distinct valuations. Continuous scaling of semivaluations
\cite[\S1.1.4]{BJ22} gives inverse continuous maps
\begin{equation}\label{eq:p081-uniform-collar}
 F\times\mathbf R_{>0}\longrightarrow\mathcal C,
 \quad(x,a)\longmapsto ax,
 \qquad y\longmapsto(a(y)^{-1}y,a(y)).
\end{equation}
The fibre here is the ordinary analytic fibre of
\cite[Definition/Exercise~4.1.7.1]{Temkin2011}: the valued field
\(k(t)\) at \(\eta\) completes to \(K\). Affine-chart seminorms
over \(\eta\) localize over \(k(t)\) and extend to that completion,
giving the charts of \(X_K^{\mathrm{an}}\); restriction reverses the
construction. The topology is the topology of evaluations of regular
functions, hence is the usual fibre subspace topology. No Gauss-section
image is identified with this full fibre.

\emph{Why usc regularization commutes for this controlled family.}
Let \(f_i\in[-M,0]\) be FS and choose the extensions \(\Phi_i\)
above with one \(C\). Every normalized section weight \(\beta_{ij}\)
has absolute value at most \(C\). The trivial metric is homogeneous:
\(\log|s|_0(ax)=a\log|s|_0(x)\), by evaluation at the unchanged
centre; see \cite[Remark~1.24 and \S1.1.4]{BJ22}. Comparing the finite
maxima with weights \(\beta_{ij}\) and \(a\beta_{ij}\) gives
\[
 |\Phi_i(ax)-a\Phi_i(x)|\leq C|a-1|,
 \qquad |\Phi_i(ax)-\Phi_i(x)|\leq2C|a-1|\quad(x\in F).
\]
The second inequality uses \(\Phi_i(x)=f_i(x)\in[-M,0]\).
Terms equal to \(-\infty\) do not affect the estimate; some term is
finite because the sections generate. The bound is uniform in \(i\).

Put \(V=\sup_i\Phi_i\), \(g=(V|_F)^{*F}\). Taking suprema gives
\(|V(ax)-V(x)|\leq2C|a-1|\). At \(x\in F\), choose a fibre
neighbourhood \(W\) with \(V|_F\leq g(x)+\epsilon\). On the collar
neighbourhood \(W\times\{|a-1|<\epsilon/(2C)\}\) one has
\(V\leq g(x)+2\epsilon\). Thus \(V^{*S}(x)\leq g(x)\). The
reverse inequality follows by restricting ambient neighbourhoods to the
fibre. This proves \eqref{eq:p081-uniform-usc-fibre} with neighbourhoods,
not sequences. For arbitrary functions only
\((V|_F)^{*F}\leq V^{*S}|_F\) follows formally. The reverse inequality
here uses the collar and the uniform weight estimate, not continuity of
the fibre inclusion alone. For example, on this collar the continuous
functions \(h_n(x,a)=\min\{1,n|a-1|\}\) vanish on \(F\), whereas
their supremum is one off \(F\). Its ambient usc regularization is one
on \(F\), while its fibre regularization is zero. This is only a
topological example, not an FS or psh family in the fixed class.

\emph{Bounded FS families and explicit psh restriction.}
By the assumed EP on \(S\),
\(\Psi=V^{*S}\in\operatorname{PSH}(S,c_1(A))\), with
\(-C\leq\Psi\leq1\). Since \(A\) is ample,
\cite[Theorem~4.15(iii) and Example~3.9]{BJ22} give a decreasing net
\(\Psi_j\in H^{\mathrm{dom}}(A)\subset H_{\mathbf Q}(A)\) with limit
\(\Psi\). Each restricts to an FS potential on the fixed \(L_K\): its
globally generating section system remains generating on the fibre, and
the frame \(U^{mC}\) for the base factor has norm one there. Thus the
restriction is exactly the FS formula with the constant-model reference.
The restricted net decreases to \(\Psi|_F=g\), so
\(g\in\operatorname{PSH}_{\mathrm{FS}}(L_K)\). No pre-existing psh
restriction or envelope theorem over \(K\) has been used.

\emph{Elementary closure of the target class.}
It is closed under further decreasing limits when the limit is not
identically \(-\infty\): apply the abstract compact-space diagonal
lemma \cite[Lemma~4.6]{BJ22} with one fixed class parameter and the
continuous FS potentials as approximants. This lemma does not assume
EP; see also the exact fixed-metric version
\cite[Lemma~7.4 and Corollary~7.5]{BE2021}. Finite maxima are treated using products of approximation nets and
concatenating section systems after common tensor powers. Constants are
FS, using a globally generated power of \(L\) with equal weights on a
generating \(k\)-system. Addition of constants is allowed. Decreasing
limits of continuous functions are usc.

\emph{Bounded psh families: a product-net argument.}
Let \(-M\leq u_i\leq0\) be a nonempty family in the target class,
and choose decreasing FS nets \(u_{i,j}\downarrow u_i\). They are
already at least \(-M\). For each \(i\), there is a cofinal tail
bounded above by one: the open sets \(\{u_{i,j}<1\}\) cover compact
\(F\); a finite subcover and directedness give one index beyond which
the bound holds everywhere. Restrict each index set \(J_i\) to this tail.

The coordinatewise ordered product \(J=\prod_iJ_i\) is directed. Put
\(W_{\mathbf j}=(\sup_i u_{i,j_i})^{*F}\).
Each is psh by the bounded-FS case after shifting by one; the resulting
FS family lies in \([-M-1,0]\). These \(W_{\mathbf j}\) decrease and
are at least \(-M\), so their limit \(W\) is psh. If
\(s=\sup_i u_i\), then \(W\geq s^{*F}\).

For the converse, let \(v\) be any continuous majorant of \(s\) and
let \(\epsilon>0\). For each \(i\), the open sets
\(\{u_{i,j}<v+\epsilon\}\) cover \(F\). Compactness and directedness
give \(j_i\), within the selected tail, such that
\(u_{i,j_i}\leq v+\epsilon\) everywhere. Choose all these indices.
Then \(W_{\mathbf j}\leq v+\epsilon\), since the right side is
continuous. Thus \(W\leq v\).

For a bounded function \(s\) on compact Hausdorff \(F\),
\begin{equation}\label{eq:p081-uniform-majorants}
 s^*=\inf\{v\in C^0(F):v\geq s\}.
\end{equation}
Indeed, every such \(v\) majorizes \(s^*\). At a fixed point \(x\),
upper semicontinuity gives a neighbourhood where
\(s^*<s^*(x)+\delta\). A continuous bump equal to zero at \(x\) and
one outside this neighbourhood exists by normality. Add a sufficiently
large multiple of it to \(s^*(x)+\delta\), obtaining a continuous
majorant with that value at \(x\). Let \(\delta\) tend to zero.
Equation \eqref{eq:p081-uniform-majorants} now gives \(W=s^{*F}\).
No countable cofinal subset or unjustified interchange of infimum and
supremum was used; coordinate choices use ordinary set-theoretic choice.

\emph{Remove the lower bound.}
For a general family bounded above, subtract its common upper bound.
The clamps \(\max\{u_i,-M\}\), \(M\in\mathbf Z_{>0}\), are psh
and lie in \([-M,0]\). Consequently
\[
 (\sup_i\max\{u_i,-M\})^{*F}
   =\max\{(\sup_i u_i)^{*F},-M\}
\]
is psh. The identity follows directly from the neighbourhood definition
of usc regularization, also at values \(-\infty\). These functions
decrease to \((\sup_i u_i)^{*F}\), which is not identically
\(-\infty\) because the family is nonempty and it majorizes any one
member. Decreasing-net closure concludes the proof. Continuous-obstacle
continuity, arbitrary fibrewise psh lifting, arbitrary-field comparison,
and a reconstruction from Gauss traces are not inputs.
\end{proof}
